\section*{Chapter \ref{ch:b1:elementary-steps} --- Reaction Mechanisms: Elementary Steps and Approximations}

\begin{solution}{exo:b1:elementary-steps:1}
(a) Unimolecular: can be elementary. (b) Bimolecular: can be elementary.
(c) Three molecules and four bonds rebuilt at once: not elementary.
(d) Termolecular, \ce{N2} carrying away the energy released: a rare but
genuine \omterm{def:b1:elementary-steps:elementary-step}{elementary step}.
\end{solution}

\begin{solution}{exo:b1:elementary-steps:2}
$v = k[\mathrm{A}][\mathrm{B}]$; $v = k[\mathrm{A}]^2$; $v = k[\mathrm{A}]$.
An overall equation is a balance sheet, not a description of how the
molecules meet; its \omterm{def:b1:rate-laws:order}{rate law} must be measured, or derived from a
mechanism.
\end{solution}

\begin{solution}{exo:b1:elementary-steps:3}
The intermediate is the hollow at 30; the \omterm{def:b1:elementary-steps:energy-profile}{transition states} are the
maxima at 70 and 55. The highest \omterm{def:b1:elementary-steps:energy-profile}{transition state} is the first: the
first step is rate-determining. It is endothermic (from 0 up to 30).
\end{solution}

\begin{solution}{exo:b1:elementary-steps:4}
By the same factor $10^4$ (\cref{prop:b1:elementary-steps:catalysis}); the
equilibrium constant is unchanged; equilibrium is reached about $10^4$
times sooner.
\end{solution}

\begin{solution}{exo:b1:elementary-steps:5}
$t_{\max} = \ln(0.30/0.10)/(0.30 - 0.10) = \qty{5.49}{s}$. Then
$[\mathrm{B}]_{\max} = a_0\frac{0.10}{0.20}(\eu^{-0.549} - \eu^{-1.648}) =
0.5\,a_0(0.577 - 0.192) = 0.192\,a_0$.
\end{solution}

\begin{solution}{exo:b1:elementary-steps:6}
$\dd[\mathrm{I}]/\dd t = k_1[\mathrm{A}] - (k_{-1} + k_2)[\mathrm{I}] = 0$
gives $[\mathrm{I}] = k_1[\mathrm{A}]/(k_{-1} + k_2)$ and $v = k_2[\mathrm{I}]
= \frac{k_1k_2}{k_{-1}+k_2}[\mathrm{A}]$. If $k_2 \gg k_{-1}$, $v =
k_1[\mathrm{A}]$: the first step is rate-determining. If $k_2 \ll k_{-1}$,
$v = (k_1/k_{-1})k_2[\mathrm{A}]$: a pre-equilibrium followed by a slow
step.
\end{solution}

\begin{solution}{exo:b1:elementary-steps:7}
Sum: \ce{2H2O2 -> 2H2O + O2}. \omterm{def:b1:elementary-steps:catalyst}{Catalyst}: \ce{I-} (consumed, then
regenerated); intermediate: \ce{IO-}. The slow first step fixes the
rate: $v = k_1[\ce{H2O2}][\ce{I-}]$.
\end{solution}

\begin{solution}{exo:b1:elementary-steps:8}
The step is strongly uphill: by the Hammond postulate its \omterm{def:b1:elementary-steps:energy-profile}{transition
state} resembles the carbocation. A substituent that lowers the energy of
the carbocation lowers that of the \omterm{def:b1:elementary-steps:energy-profile}{transition state} almost as much, so
the barrier falls and the step is faster.
\end{solution}

\begin{solution}{exo:b1:elementary-steps:9}
$k_{\mathrm B}/k_{\mathrm C} = \exp(20000/RT)$: \num{3.0e3} at \qty{300}{K},
$\exp(4.01) = 55$ at \qty{600}{K}. Heating reduces the kinetic preference
for B and, by making the formation of B reversible, lets the system reach
the more stable C.
\end{solution}

\begin{solution}{exo:b1:elementary-steps:10}
$k_1[\mathrm{A}][\mathrm{M}] = (k_{-1}[\mathrm{M}] + k_2)[\mathrm{A}^*]$,
so $v = k_2[\mathrm{A}^*] = \dfrac{k_1k_2[\mathrm{A}][\mathrm{M}]}{k_{-1}[\mathrm{M}]
+ k_2}$. At high pressure ($k_{-1}[\mathrm{M}] \gg k_2$), $v =
(k_1k_2/k_{-1})[\mathrm{A}]$: order 1. At low pressure, $v =
k_1[\mathrm{A}][\mathrm{M}]$: order 2, activation by collision becomes the
slow step.
\end{solution}

\begin{solution}{exo:b1:elementary-steps:11}
The proof is that of the theorem. For $k_1 = k_2 = k$:
$\frac{\dd}{\dd t}([\mathrm{B}]\eu^{kt}) = ka_0$, so $[\mathrm{B}]\eu^{kt} =
ka_0t$ and $[\mathrm{B}] = a_0kt\,\eu^{-kt}$, maximal at $t = 1/k$.
\end{solution}

\begin{solution}{exo:b1:elementary-steps:12}
$\mathrm{A} + \mathrm{B} \rightleftharpoons \mathrm{I} + \mathrm{C}$ (fast,
constant $K_1$), then $\mathrm{I} + \mathrm{B} \to \mathrm{P}$ (slow,
$k_2$). The sum is the overall reaction; $[\mathrm{I}] =
K_1[\mathrm{A}][\mathrm{B}]/[\mathrm{C}]$ and $v = k_2[\mathrm{I}][\mathrm{B}]
= k_2K_1[\mathrm{A}][\mathrm{B}]^2/[\mathrm{C}]$: the by-product, released
in the pre-equilibrium, slows the reaction.
\end{solution}

\begin{solution}{pb:b1:elementary-steps:1}
\textbf{1.} Doubling $[\ce{NO}]_0$ multiplies $v_0$ by 4: order 2; doubling
$[\ce{O2}]_0$ multiplies it by 2: order 1.
\textbf{2.} $v = k[\ce{NO}]^2[\ce{O2}]$; $k = 1.0 \times 10^{-5}/((1.0
\times 10^{-3})^2 \times 1.0 \times 10^{-3}) =
\qty{1.0e4}{L^2.mol^{-2}.s^{-1}}$.
\textbf{3.} A simultaneous collision of three molecules is rare, and a
single step could not get slower on heating (question 5).
\textbf{4.} $E_a = R\ln(k_{400}/k_{300})/(1/300 - 1/400) = 8.314 \ln(0.1)
/(8.33 \times 10^{-4}) = \qty{-23}{kJ/mol}$.
\textbf{5.} An \omterm{def:b1:elementary-steps:elementary-step}{elementary step} proceeds through collisions that cross a
barrier; the fraction of them able to do so, $\eu^{-E_a/RT}$ with $E_a \ge
0$, can only grow with temperature.
\textbf{6.} $-\dd[\ce{NO}]/\dd t = 2v$, $-\dd[\ce{O2}]/\dd t = v$.
\textbf{7.} \ce{2NO -> N2O2} plus \ce{N2O2 + O2 -> 2NO2} gives
\ce{2NO + O2 -> 2NO2}; the intermediate is \ce{N2O2}.
\textbf{8.} Both bimolecular (and the reverse of the first unimolecular).
\textbf{9.} $v = k_2[\ce{N2O2}][\ce{O2}]$.
\textbf{10.} $[\ce{N2O2}] = K_1[\ce{NO}]^2$ with $K_1 = k_1/k_{-1}$.
\textbf{11.} $v = k_2K_1[\ce{NO}]^2[\ce{O2}]$: the observed law.
\textbf{12.} $k = k_1k_2/k_{-1}$.
\textbf{13.} $\dd[\ce{N2O2}]/\dd t = k_1[\ce{NO}]^2 - k_{-1}[\ce{N2O2}] -
k_2[\ce{N2O2}][\ce{O2}]$.
\textbf{14.} Setting it to zero: $[\ce{N2O2}] = k_1[\ce{NO}]^2/(k_{-1} +
k_2[\ce{O2}])$.
\textbf{15.} $v = k_2[\ce{N2O2}][\ce{O2}]$ gives the stated law.
\textbf{16.} $k_{-1} \gg k_2[\ce{O2}]$: the dimer falls apart far more often
than it meets dioxygen, so the first step stays at equilibrium.
\textbf{17.} If $k_2[\ce{O2}] \gg k_{-1}$, $v = k_1[\ce{NO}]^2$: order 2 in
\ce{NO} and 0 in \ce{O2}.
\textbf{18.} The measured order 1 in \ce{O2}: the first limit.
\textbf{19.} $\ln k = \ln k_1 + \ln k_2 - \ln k_{-1}$.
\textbf{20.} Differentiating with respect to $T$ and using $\dd\ln k_i/\dd T
= E_i/RT^2$: $E_a = E_1 + E_2 - E_{-1}$.
\textbf{21.} Heating speeds the slow step a little ($E_2$ small) but shifts
the pre-equilibrium strongly back towards \ce{NO} ($E_{-1}$ large): fewer
dimers, and the net rate falls.
\textbf{22.} The dimer lies $E_1 - E_{-1} = \qty{-38}{kJ/mol}$ below the
reactants and the second \omterm{def:b1:elementary-steps:energy-profile}{transition state} $E_2 = \qty{15}{kJ/mol}$ above
the dimer: \qty{23}{kJ/mol} below the reactants. The highest point of the
path is the first \omterm{def:b1:elementary-steps:energy-profile}{transition state}, only \qty{2}{kJ/mol} up.
\textbf{23.} $\exp\big(\frac{23000}{8.314}(\frac{1}{400} - \frac{1}{300})\big)
= 0.10$: ten times slower at \qty{400}{K}.
\textbf{24.} $E_a = 2 - 40 + 15 = \textbf{\qty{-23}{kJ/mol}}$, the value
measured in Part I: the mechanism explains both the \omterm{def:b1:rate-laws:order}{rate law} and its
strange temperature dependence.
\end{solution}
